\section{Scope and Coverage}

Every theory has edges, and a reader is owed an honest map of where they fall. This section draws that map. It states plainly what the model establishes, what it extrapolates, and what it does not touch at all, so that no one mistakes silence for a claim, and no one reads a demonstration of shape as a quantitative fit it was never meant to be.

The thesis of the whole work is a thesis about origin, not about reproduction. The claim is that \textbf{one discrete substrate}, the vibe mesh, a flow of the eight atoms, can give rise to the broad \textbf{shapes} of known physics, and, more tentatively, to the shapes of life and mind. The substrate is the $\{3,4,3,4\}$ hyperbolic honeycomb, its docks carrying a ternary tone in $\{-1,0,+1\}$, evolving under the knit, the collide-then-stream law, with the mesh growing at its wake. From those few ingredients the right \emph{kinds} of structure appear. The model does not attempt a system-by-system numerical reproduction of the established sciences. It is a claim about the family a world belongs to, not a recomputation of every member.

\begin{principle}[Breadth, not depth]
The claim is breadth of origin, not depth of reproduction. The model aims to show that the right kinds of structure across physics, and more tentatively across life and mind, follow from one substrate. It does not claim to compute chemistry, atomic spectra, or any absolute physical constant.
\end{principle}

\subsection{What is established}

The first tier holds the results that are either derived analytically on the $\{3,4,3,4\}$ substrate or measured directly in simulation on it. These are not extrapolations. They are the structural skeleton of fundamental physics, read off the base.

\textbf{Spin and matter.} The $24$ sites of a dock form the $D_4$ root system, the $24$-cell, which decomposes under $\mathrm{SO}(8)$ as $8v + 8s + 8c$. The two spinor eights are genuine half-integer-spin representations, so the substrate carries spinors at the base rather than assuming them. Triality, the order-three outer symmetry permuting the three eights, seeds three generations. The self appears here as a topological soliton of odd charge, hence a fermion by spin-statistics, and a Skyrme-type stabilizer term is measured positive on the substrate, so matter is stable rather than collapsing.

\textbf{Causality and dynamics.} A \textbf{finite light cone} with dynamical exponent $z = 1$, and the Dirac dispersion that goes with it, are read off the knit directly. Signals propagate at one dock per beat, and the linearized response around the vacuum reproduces the relativistic dispersion relation.

\textbf{Unification and mass.} The substrate exhibits a \textbf{grand unification} breaking to the Standard Model gauge content, with the weak mixing angle fixed at $\sin^2\theta_W = 3/8$ at the unification scale. The $b$-$\tau$ mass relation and the determinant mass \textbf{ratios} follow, and three generations descend from triality. These are ratios and scale relations, not absolute masses.

\textbf{Gravity and force.} A $1/r$ Newtonian gravity emerges in the flat cubic cusp, while a $1/r^2$ holographic correlator governs the negative-curvature bulk. The model delivers $\mathrm{U}(1)$ electromagnetism, $D_4$ isotropy of the cusp, de Sitter expansion from the growth of the wake, the holographic hierarchy, and a resolution of the Bell tension at the boundary.

\textbf{Substrate selection.} The choice of $\{3,4,3,4\}$ itself is established, not assumed. Three independent constraints converge on it. The dimension ladder, the crystallographic requirement, and the spinor requirement together single out this honeycomb among the alternatives.

\begin{result}[The established skeleton]
On the $\{3,4,3,4\}$ substrate the model derives or measures spinors and three generations, a finite light cone with the Dirac dispersion, grand unification with $\sin^2\theta_W = 3/8$ and the mass ratios, $1/r$ gravity in the cusp with a $1/r^2$ correlator in the bulk, electromagnetism, isotropy, de Sitter growth, the holographic hierarchy, and the substrate selection itself.
\end{result}

\subsection{What is extrapolated}

The second tier holds results that are plausible and whose ingredients are demonstrated, but whose full chain is not run end to end. These are grounded extrapolation, not proof, and they are labeled as such.

The principal extrapolation is the emergent ladder that climbs above physics. The full ascent runs from atoms to molecules to cells to tissues to senses to minds. Only the first two rungs are simulated end to end on the substrate. The field-to-particle step is run, and then the particle-to-composite step is run. The higher rungs each have their \textbf{ingredient} demonstrated in isolation. Binding is shown. Self-maintenance is shown. Copying is shown. The onion of senses is shown, and the integrated hub is shown. What is not done is running them together as one continuous simulation from tone to mind. The ladder is assembled from demonstrated parts, not evolved as a single object.

\begin{principle}[Grounded extrapolation]
Life, mind, intelligence, and consciousness appear in the model as regimes of the one substrate, demonstrated at toy scale and to first order. Each rung's mechanism is shown. The unbroken chain from the ternary tone up to a mind is not run as one simulation. These results are grounded extrapolation, not proof.
\end{principle}

\subsection{What is not covered}

The third tier is the honest list of what the model does not touch, stated so that no reader mistakes an omission for a hidden claim. The self studied here is the minimal conserved knot, not hydrogen, and the model makes no pretense of computing the real periodic table.

\begin{table}[ht]
\centering
\begin{tabular}{@{}p{0.32\textwidth}p{0.6\textwidth}@{}}
\toprule
area & what is out of scope \\
\midrule
atomic and nuclear structure & no real atoms, no binding energies, no shell model, no periodic table. The self here is the minimal conserved knot, not hydrogen \\
chemistry & no molecular orbitals, bonds, reaction rates, or thermodynamics of real substances \\
condensed matter & no band structures, phases, or transport coefficients of real solids \\
quantitative particle data & no absolute masses, no full CKM or PMNS mixing, no CP-violation magnitudes, no decay rates. Only ratios and the unification-scale relations \\
full numerical relativity & no binary merger or forming horizon evolved dock-by-dock from the tone \\
quantitative cosmology & no radiation or matter eras, no nucleosynthesis, no CMB spectrum, no structure formation. Only the de Sitter growth law \\
absolute scales & the beat-to-seconds and dock-to-meters identifications are not fixed, so all rates stay in beat units \\
real biology and neuroscience & the life and mind results are mechanism demonstrations at toy scale, not models of any actual organism, cell, or brain \\
the exact knit's outputs & the precise Standard Model beta coefficients, absolute masses, and the settling of matter into cusps stay open \\
\bottomrule
\end{tabular}
\caption{What the model does not cover.}
\label{tab:out-of-scope}
\end{table}

A word on general relativity is needed, because its status is easy to misread. The Einstein field equations \emph{as an equation of state}, the bending of light, black-hole thermodynamics, and gravitational-wave forms \textbf{are} derived on the substrate. So general relativity is in scope at the effective-theory level, and its quantitative relations are recovered. What is not done, and what no one does from a micro-law, is evolving a binary merger or a forming horizon dock-by-dock directly from the ternary tone. The effective theory is established. The dock-by-dock evolution is the frontier, not a gap in the claim.

Two items once listed as scope limits are now resolved, so they are removed from the list. The $\{3,4,3,4\}$ addressing is solved and verified, so the geometry is fully navigable in simulation. And the substrate is shown computationally universal on its own ternary knit, so the law is known to be capable of arbitrary computation without any added machinery.

\subsection{The three-layer robustness audit}

Not every established result leans on the same assumptions, and the reader deserves to know exactly how much each one depends on. Every result sits in one of \textbf{three layers}, sorted by what it rests on. The layer is the measure of robustness. A bulk result survives almost any choice. A flat-layer result needs the cusp to be clean.

\begin{table}[ht]
\centering
\begin{tabular}{@{}p{0.14\textwidth}p{0.28\textwidth}p{0.5\textwidth}@{}}
\toprule
layer & depends on & what lives here \\
\midrule
\textbf{bulk} & any hyperbolic substrate & spin, triality, the $\mathrm{SO}(10)$ group, $\sin^2\theta_W = 3/8$, the mass ratios, the $1/r^2$ holographic correlator, cosmology, the hierarchy, the toolkit \\
\textbf{boundary} & the negative-curvature boundary & holography, the resolution of the Bell tension \\
\textbf{flat} & the flat cubic cusp being clean periodic 3D space & the Dirac equation, the light cone, Lorentz invariance, $1/r$ gravity, electromagnetism, solitons, self-as-fermion, matter formation \\
\bottomrule
\end{tabular}
\caption{The three robustness layers, sorted by what each result depends on.}
\label{tab:robustness-layers}
\end{table}

\begin{principle}[Only the flat layer leans on the cusp]
Only the flat-layer results lean on the cubic cusp being clean periodic three-space. The bulk and boundary results do not reach the cusp at all. They follow from the hyperbolic geometry and its boundary alone, so they are robust against any reasonable change in the fine structure of the slice.
\end{principle}

This layering is also why the flat-layer results are a measured fit and not a tuned one. The flat-layer physics needs the cusp to look like clean periodic three-dimensional space, and that is a property one can measure rather than assume. A clean cubic cusp gives a spectral dimension near \textbf{3} and clean \textbf{isotropy}. A generic disordered slice does not, and the flat-layer physics visibly degrades on it.

\begin{table}[ht]
\centering
\begin{tabular}{@{}lll@{}}
\toprule
observable & clean cubic cusp & generic disordered slice \\
\midrule
spectral dimension & about $2.97$ (clean 3D) & about $2.16$ (sub-3D) \\
light-cone isotropy & about $0.14$ (smooth) & about $0.32$ (rough) \\
\bottomrule
\end{tabular}
\caption{The flat-layer physics holds on the clean cusp and degrades on a disordered slice.}
\label{tab:cusp-vs-disordered}
\end{table}

The right surface is therefore forced rather than chosen. Only the clean periodic cusp gives the spectral dimension and isotropy that the flat-layer physics requires, and a generic slice fails both tests at once. The flat-layer results are a measured fit to a forced surface, which is what separates them from fine-tuning. The substrate is not adjusted to fit the physics. The one surface that carries the physics is the one the geometry already singles out.

\subsection{How to read the scope}

Read together, the three tiers say one thing. One substrate appears to give the right \textbf{kinds} of structure across physics, and more tentatively across life and mind, from almost nothing. It does not yet, and does not here claim to, compute chemistry or atomic spectra or any absolute constant of nature.

\begin{principle}[The honest summary]
We are showing the fire can take the shape of a world. We are not yet recreating every flame in it.
\end{principle}

\subsection{See also}

The named frontiers, the testable edges, and the consequences of the model are taken up in the sections that follow. Section~70, Open Problems, lists the frontiers by name. Section~68, Predictions and Falsifiability, gives the testable edges where the model could be broken. Section~67, Implications, states what follows if the model holds.
